% Direct-authority diplomatic French for printed p.210.
\clearpage
\SourcePage{210}

\LayerTarget{tohoku:proposition:5.5.4}
\noindent\textit{où $C({\mathbf U},P)$ désigne le complexe (à opérateurs) des
cochaines quelconques (non nécessairement alternées) de ${\mathbf U}$ à
coéfficients dans $P$.}

Conjuguant avec le corollaire 1 du théorème 5.5.3, il vient :

\LayerTarget{tohoku:corollary:fixed-point-free-g-covering}
\textsc{Corollaire.} \textit{${\mathbf U}$ étant un $G$-recouvrement
sans point fixe, on suppose que ${\mathbf U}$ est ouvert, ou que ${\mathbf U}$
est fermé et de plus $X$ paracompact ou ${\mathbf U}$ fini. Si $A$ est un
$G$-faisceau tel que $H^n(U_{i_0\ldots i_p},A)=0$ pour tout
$(i_0\ldots\,.\,i_p)\in I^{p+1}$ et tout $n>0$, alors on a des isomorphismes
canoniques :}

\LayerTarget{tohoku:equation:fixed-point-free-g-covering-cohomology}
\[
H^n(X;G,A)=H^n\bigl(C({\mathbf U},A)^G\bigr).
\]

\LayerTarget{tohoku:calculation:nonhomogeneous-equivariant-cochains}
Nous allons expliciter le complexe $C({\mathbf U},P)^G$ quand ${\mathbf U}$ est
un $G$-recouvrement sans point fixe. Alors l'ensemble d'indices de
${\mathbf U}$ est isomorphe (en tant qu'ensemble où opère $G$) à un produit
$G\times I$ où $G$ opère par translations gauches sur le premier facteur :
$g'\cdot(g,i)=(g'g,i)$; nous identifierons $i$ à $(e,i)$. Soit
$(f_{g_0i_0,\ldots,g_ni_n})$ une $n$-cochaine invariante à coéfficients dans
$P$, posons

\LayerTarget{tohoku:equation:5.5.9}
\[
\tag{5.5,9}
F_{i_0,g_1i_1,\ldots,g_ni_n}
=f_{e.\,i_0,g_1i_1,\ldots,g_ni_n}
\]
$F$ est une fonction qui dépend des $n+1$ arguments $i_0,\ldots,i_n\in I$ et
des $n$ arguments $g_1,\ldots,g_n\in G$, à valeurs dans les
$P(U_{i_0}\cap g_1U_{i_1}\cap\ldots\,.\,\cap g_{i_n}U_{i_n})$, et la cochaine
invariante donnée est entièrement déterminée par la connaissance de cette
``\emph{cochaine non homogène}'' $F$, en vertu de

\LayerTarget{tohoku:equation:5.5.10}
\[
\tag{5.5.10}
f_{g_0i_0,\ldots,g_ni_n}
=g_0.\,F_{i_0,g_0^{-1}g_1i_1,\ldots,g_0^{-1}g_ni_n}
\]
cette formule définit d'ailleurs, quand on part d'une cochaine non homogène
$F$, une cochaine invariante $f$, et la cochaine non homogène associée à $f$
n'est autre que $F$. Reste à exprimer l'opérateur différentiel de
$C({\mathbf U},P)^G$ directement en termes de cochaines non homogènes. On trouve
immédiatement :

\LayerTarget{tohoku:equation:5.5.11}
\[
\tag{5.5.11}
\begin{aligned}
&(\partial F)_{i_0\,g_1i_1,\ldots,g_{n+1}i_{n+1}}
=g_1.\,F_{i_1,g_1^{-1}g_2i_2,\ldots,
g_1^{-1}g_{n+1}i_{n+1}}\\
&\qquad+\sum_{1\leqq\alpha\leqq n+1}(-1)^\alpha
F_{i_0,g_1i_1,\ldots,\widehat{g_\alpha i_\alpha},\ldots,
g_{n+1}i_{n+1}}
\end{aligned}
\]
oû le signe $\wedge$ signifie que les lettres
placées en dessous doivent être omises.
\LayerTarget{tohoku:special-case:single-orbit-covering}
Un cas particulièrement intéressant est celui oû $I$ est réduit à un élément,
c'est à dire où on part d'une partie $D$ de $X$ telle que
$\displaystyle\bigcup_{g\in G}g\cdot D=X$, et qu'on considère le recouvrement
$(g\cdot D)_{g\in G}$. Alors les $n$-cochaines non homogènes sont des systêmes
$(F_{g_1,\ldots,g_n})$, avec

\LayerTarget{tohoku:equation:single-orbit-cochain-domain}
\[
F_{g_1,\ldots,g_n}\in
P(D\cap g_1\cdot D\cap\ldots\,.\,\cap g_n\cdot D)
\]
et la formule (5.5.11) devient

\LayerTarget{tohoku:equation:5.5.11bis}
\[
\tag{5.5.11 bis}
\begin{aligned}
&(\partial F)_{g_1,\ldots,g_{n+1}}
=g_1\cdot F_{g_1^{-1}g_2,\ldots,g_1^{-1}g_{n+1}}\\
&\qquad+\sum_{1\leqq\alpha\leqq n+1}(-1)^\alpha
F_{g_1,\ldots,\widehat{g_\alpha},\ldots,g_{n+1}}.
\end{aligned}
\]

\LayerTarget{tohoku:construction:complex-from-local-composition-data}
Soit donc $L$ l'ensemble des $g\in G$ tels que $g\cdot D\cap D\ne\phi$ (ce
sera un ensemble fini dans les cas ``raisonnables''), la formule précédente
montre que le complexe $C({\mathbf U},P)^G$ donnant les groupes
$H^n({\mathbf U};G,P)$ peut se construire en connaissant

\TohokuContinuation{211}
